\section{The core process and exact waiting-event elimination}
\label{model:section}

\subsection{States, genealogy, and the imported exponent}

Each real vertex has an independent rate-one source clock. At a ring it
chooses a neighbour uniformly and creates a new leaf at that neighbour.
The total birth rate of a full tree of size $n$ is $n$, so its embedded
arrival chain is the discrete random friend tree. Every finite tree is
reachable from an initial edge with positive probability in a fixed
finite time. Consequently the expected core count from any fixed
finite seed has upper exponential rate at most the rate from the edge.

We use the following established input from~\cite{Uzay}: for the
process started from an edge,
\begin{equation}\label{model:imported-exponent}
 F_2(t):=\mathbb E_{\mathrm{edge}}C(t)=e^{\mu t+o(t)},
 \qquad
 \mathbb E C_n=n^{\mu+o(1)},
 \qquad 0.5650\le\mu\le0.5665.
\end{equation}
The proof source is fixed at the revision in the bibliography. The
weighted planted-branch comparison used below is also stated explicitly
with its location in that source. We do not use an unweighted ordering
between planted and free branch means.

A core state consists of a finite tree $H$ of real cores, leaf counts
$L_i\ge0$, and inert-stub counts $b_i^\partial\ge0$. Put
\[
 k=|H|,\qquad c_i=\deg_H(i)+b_i^\partial,
 \qquad D_i=L_i+c_i\ge2.
\]
An inert stub contributes to $D_i$ but has no clock. A choice of that
stub produces a null transition. For a full tree all stub counts are
zero. In that case a singleton core has at least two leaves. In the
one-stub retained newborn, $k=1$, $c_1=L_1=1$, so the singleton has
exactly one real leaf. These two cases have different core-ring
transition laws and are kept separate.

Whenever a core chooses one of its leaves, that leaf becomes a new
core with one leaf of its own. Call it a core child of the selecting
core. Other arrivals only add a leaf to an existing core. This defines
a permanent genealogy of cores. A newborn tag has one leaf and one
core neighbour, either a real parent or an inert stub. Let $\mathcal X$
be the countable class of all finite environments with such a tagged
newborn, allowing finitely many inert stubs away from the tag as well.
Only the tag and its future genealogical descendants are counted in
$F_x(t)$.

For $s>0$, define a positive operator on bounded functions on
$\mathcal X$ by
\begin{equation}\label{model:tagged-kernel}
 (\mathcal K_s f)(x)
 =\mathbb E_x\sum_{v:\,v\text{ a future child of the tag}}
             e^{-s\sigma_v} f(X_v),
\end{equation}
where $X_v$ is the entire environment at the birth of $v$, retagged
at $v$. The strong Markov property and Tonelli's theorem give
\[
 (\mathcal K_s^m\mathbf1)(x)
 =\mathbb E_x\sum_{\operatorname{depth}(v)=m}e^{-s\sigma_v}
 =A_m(s;x).
\]
The whole environment is retained in $X_v$; this is not an assertion
that descendant families evolve independently.

Give every potential core a rate-one clock and let every ring create a
potential child, whether or not the corresponding actual ring promotes
a leaf. Assign unused potential branches fresh independent drivers.
Every actual tagged genealogy is a subset of this potential genealogy,
whose expected depth-$j$ count at time $t$ is $t^j/j!$. In particular
\begin{equation}\label{model:kernel-bounds}
 \|\mathcal K_s\|_\infty\le1/s,
 \qquad [2(s+2)]^{-m}\le A_m(s;x)\le s^{-m}.
\end{equation}
For the lower bound, wait for the first ring of the tag or its initial
leaf, and require a tag ring choosing a leaf. Before that ring the
tag still has one core neighbour, while outside arrivals can only add
leaves. Its leaf-choice probability is at least $1/2$. The discounted
probability is therefore at least $1/[2(s+2)]$. Iterate at newborn
states. Core-count Yule domination also gives $F_x(t)\le e^t$ uniformly.

Finite inert environments introduce no larger fixed-time supremum.
Replace each of their $b$ stubs by an edge to a genuine auxiliary
parent of degree $R$, completed by $R-1$ leaves. Until a return across
one of these edges, the retained histories agree. The probability of
such a return by a fixed time $T$ is at most $bT/R$, since degrees
never decrease. Uniform tagged Yule second moments and
Cauchy--Schwarz make the difference in tagged expected rewards tend
to zero as $R\to\infty$. Thus every fixed inert profile is a
finite-horizon limit of genuine finite environments.

\subsection{The integrated leaf burst}

Between two core-clock rings, the core skeleton is fixed. Each leaf
creates leaves at its own core at rate one, so the coordinates $L_i$
are independent Yule processes. The next core ring has an independent
exponential waiting time of rate $k$. Write $\mathcal Y_t$ for the
leaf semigroup and set
\begin{equation}\label{cc:R-def}
 (R_{H,s}f)(L)=\int_0^\infty e^{-(k+s)t}(\mathcal Y_t f)(L)\,dt,
 \qquad \alpha=k+s,\quad \ell=\sum_iL_i.
\end{equation}
If $I=\{i:L_i>0\}$ and $\ell>0$, the exact matrix entry is
\begin{equation}\label{cc:beta-entry}
 R_{H,s}(L,M)=
 \left[\prod_{i\in I}\binom{M_i-1}{L_i-1}\right]
 \mathrm B(\alpha+\ell,|M|-\ell+1)
\end{equation}
when $M_i\ge L_i$ for $i\in I$ and $M_i=0$ otherwise, and it is zero
off this support. If $\ell=0$, the leaf vector remains zero and the
only entry is $R_{H,s}(0,0)=1/\alpha$.

Indeed, with $z=e^{-t}$, the Yule transition probability from $L_i>0$
to $M_i$ is
$\binom{M_i-1}{L_i-1}z^{L_i}(1-z)^{M_i-L_i}$.
Multiplication and integration against $dt=-dz/z$ give
\eqref{cc:beta-entry}. Thus leaf-clock waiting events are integrated
exactly, including their unbounded population growth.

At rational $s$, every entry is rational, since
\[
 \mathrm B(a,b+1)=\frac{b!}{a(a+1)\cdots(a+b)}
 \qquad(a>0,\ b\in\mathbb Z_{\ge0}).
\]
The exact total-leaf tail, for $N\ge\ell>0$, is
\begin{equation}\label{cc:exact-tail}
 \alpha\sum_{|M|\ge N}R_{H,s}(L,M)
 =\prod_{j=\ell}^{N-1}\frac{j}{j+\alpha}.
\end{equation}
To see this, observe the total-leaf Yule process at an independent
rate-$\alpha$ exponential time. Its hitting time of $N$ is the sum
of independent exponentials of rates $\ell,\ldots,N-1$.

\subsection{The exact core-ring kernel}

At a core ring with source $j$, every real core neighbour $i$ receives
a leaf with probability $1/(M_j+c_j)$; an inert choice is null; and
a promotion occurs with probability $M_j/(M_j+c_j)$. Define the rate
operators
\begin{align}
 (A_Hf)(M)
 &=\sum_j\left[
   \sum_{i\sim_Hj}\frac{f(M+e_i)}{M_j+c_j}
   +\frac{b_j^\partial f(M)}{M_j+c_j}\right],
   \label{model:A}\\
 (B_Hg)(M)
 &=\sum_{j:M_j>0}\frac{M_j}{M_j+c_j}
      g(H^{(j)},M^{(j)}).
   \label{model:B}
\end{align}
Here $H^{(j)}$ adjoins a child at $j$; $M^{(j)}$ decreases $M_j$ by one
and appends the new child's leaf count one. The stub counts are inherited.
Their row masses sum to $k$. Hence
\begin{equation}\label{model:Q}
 Q_s=R_s(A+B),\qquad Q_s\mathbf1(y)=\frac{C(y)}{C(y)+s}
\end{equation}
is the exact discounted transition kernel between consecutive core
rings. The notation $R_s$ denotes the appropriate $R_{H,s}$ in each
state. Null stub choices are included and consume one ring.

If desired, nonpromotion core rings can also be eliminated. With
$W_{H,s}=R_{H,s}A_H$ and $V_{H,s}=R_{H,s}B_H$, the exact first-promotion
kernel is
\begin{equation}\label{model:first-promotion}
 \mathcal J_{H,s}=\sum_{q\ge0}W_{H,s}^qV_{H,s}.
\end{equation}
Before the first promotion, the promotion intensity
$b_H(M)=\sum_jM_j/(M_j+c_j)$ does not decrease. Thus on the upward
cone above $L$,
\[
 W_{H,s}\mathbf1\le\frac{k-b_H(L)}{k+s}<1\qquad(s>0).
\]
The omitted mass after $q=0,\ldots,Q-1$ is at most
$[(k-b_H(L))/(k+s)]^Q$. The separate total-leaf tail
\eqref{cc:exact-tail} makes every finite calculation effective.
We use the full ring kernel $Q_s$ when genealogical marks are required;
its histories have explicit nonnegative rational weights at rational
$s$. Section~\ref{calc:section} gives a complete coefficient calculator.
